\subsection{A level-nine character ambiguity at the smallest odd ramification}

\begin{lemma}[Cyclic isogenies and a Tate comparison]
\label{lem:B-RD29-isogeny-Tate}
Let $K$ be a complete discretely valued field of characteristic zero,
let $r\geq2$, and let $E/K$ have a $K$-rational point $P$ of exact
order $r$. Put $B=E/\langle P\rangle$. There is a split Tate curve
$T_q/K$, with a parameter $q\in K^\times$ of positive valuation,
such that $B[r]\simeq T_q[r]$ as finite group schemes over $K$.
The generic $p$-divisible group of a Tate curve satisfies the
manuscript monodromy condition: its two graded pieces are
$\mu_{p^\infty}$ and $\mathbf Q_p/\mathbf Z_p$.
\end{lemma}
\begin{proof}
The dual of the quotient isogeny is a morphism
$\widehat\phi:B\to E$ with kernel canonically Cartier dual to
$\langle P\rangle\simeq\mathbf Z/r$. Its restriction gives an
exact sequence of finite etale $K$-group schemes
\[
 0\longrightarrow\mu_r\longrightarrow B[r]
   \xrightarrow{\widehat\phi}\langle P\rangle
   \longrightarrow0.                                      \tag{B-RD29.1}
\]
It is exact on the right: an inverse image of $P$ under
$\widehat\phi$ is killed by $r$, since
$[r]=\phi\widehat\phi$ on $B$ and $\phi(P)=0$.
Identify the last group with $\mathbf Z/r$ using its generator $P$.
The sequence is killed by $r$, so it splits as a sequence of
abstract $\mathbf Z/r$-modules. Its Galois descent class lies in
\[
 H^1(K,\mu_r)=K^\times/(K^\times)^r.
\]
Choose a representative $q_0$ and multiply it by an $r$th power
of a sufficiently high power of a uniformizer to obtain $v_K(q)>0$.
The finite group associated to this Kummer class is exactly
$T_q[r]$. This is the explicit finite group scheme
$G_{K,q,r}$ in \cite[``Points of finite order'', pp.18--19]{Tate1995Uniformization};
its component over $a\in\mathbf Z/r$ is given by $z^r=q^a$.
That source proves both its extension-class description and its
identification with the torsion of the Tate curve. The latter
curve exists by its Theorem~1, p.6.

For $r=p^n$ these descriptions at every level yield
\[
 0\longrightarrow\mu_{p^\infty}\longrightarrow T_q[p^\infty]
   \longrightarrow\mathbf Q_p/\mathbf Z_p\longrightarrow0.
\]
Both graded groups have good models over $\mathcal O_K$.
Their finite inertia characters are trivial, so the manuscript
character of $T_q[p^\infty]$ is constantly two on inertia.
\end{proof}

\begin{proposition}[B-RD29: nine-torsion fails already at $p=3,e=2$]
\label{prop:B-RD29-p3-level-two}
\status{Proved; independent parent complete proof, original-source and exact-certificate audits passed.}
Let $K=\mathbf Q_3(\zeta_3)$, whose absolute ramification is two
and whose residue field is $\mathbf F_3$. Define
\[
 B:\quad y^2-xy+y=x^3-x^2-28x-53,
 \qquad q=2^5\,3^6=23328,
\]
and let $T_q/K$ be the split Tate curve with this parameter.
Then
\[
 B[9]\simeq T_q[9],                                      \tag{B-RD29.2}
\]
but the manuscript characters of $B[3^\infty]$ and
$T_q[3^\infty]$ differ on an inertia element of order two in
the finite inertia image. The first group is potentially good;
the second satisfies the filtration condition with good rank-one
graded pieces. Consequently
\[
                 \nu_{\min}(3,2,h)\geq2\qquad(h\geq2).
                                                        \tag{B-RD29.3}
\]
The same lower bound holds for every $e\equiv2\pmod4$.
\end{proposition}

\begin{proof}
We give the quotient, the Kummer class and the inertia calculation
explicitly. Start with
\[
 E:\quad y^2-3xy-12y=x^3-12x^2,
 \qquad P=(0,0).
\]
Exact addition gives
\[
 \begin{array}{c|rrrrrrrr}
 a&1&2&3&4&5&6&7&8\\\hline
 x(aP)&0&12&4&6&6&4&12&0\\
 y(aP)&0&48&8&18&12&16&0&12
 \end{array}
\]
and $9P=0$. In particular $P$ has exact order nine.
The quotient by this subgroup is
\[
 E':\quad y^2-3xy-12y=x^3-12x^2-2310x-36190.
                                                        \tag{B-RD29.4}
\]
Here is a direct verification of the quotient formula, so that
no equation is imported without checking its isogeny. In short
coordinates $X=x-13/4$, $Y=y-(3x+12)/2$, put
\[
 A=-219/16,\quad C=827/32,\quad F(X)=X^3+AX+C,
 \quad\mathcal R=\{-13/4,35/4,3/4,11/4\}.
\]
Define
\[
 R(X)=X+\sum_{a\in\mathcal R}
   \left(\frac{2(3a^2+A)}{X-a}
              +\frac{4F(a)}{(X-a)^2}\right).
\]
Exact polynomial arithmetic gives
\[
 F(X)R'(X)^2=R(X)^3-\frac{37179}{16}R(X)
                         -\frac{1397493}{32}.
                                                        \tag{B-RD29.5}
\]
The reduced numerator and denominator of $R$ have degrees nine
and eight, and the denominator has double zeros precisely at
$\mathcal R$. Thus $(X,Y)\mapsto(R(X),YR'(X))$ extends to a
pointed degree-nine morphism of smooth elliptic curves, hence an
isogeny, and its kernel consists of the nine displayed points.
The target short equation is exactly (B-RD29.4) in the original
long coordinates. Finally $x=9x_1+1$, $y=27y_1+21$ changes
(B-RD29.4) to the displayed equation for $B$. This identifies
$B$ with $E/\langle P\rangle$ over $\mathbf Q$.

To compute the extension class in (B-RD29.1), consider
\[
 f(x,y)=\frac{(x^2-y)^3(y-3x+4)}{(x-4)^3}.
                                                        \tag{B-RD29.6}
\]
The tangent and chord construction gives
\[
 \operatorname{div}\left(\frac{x^2-y}{x-4}\right)
       =3[P]-[3P]-2[0],
 \qquad
 \operatorname{div}(y-3x+4)=3[3P]-3[0].
\]
For the first equality, the usual two addition lines are $y=0$
and $y=4x$, and the equation of $E$ gives
$y(y-4x)=(x-12)(x^2-y)$. For the second, substitution of
$y=3x-4$ into the equation gives $-(x-4)^3$.
Consequently $\operatorname{div}(f)=9[P]-9[0]$.

Let $\delta:E(K)\to K^\times/(K^\times)^9$ be the connecting
homomorphism of the dual isogeny. The precise divisor description
of this map is \cite[Appendix~A.1, Lemmas~A.1--A.2, p.22]{AntieauAuel2021Descent}.
There is a constant $c$ such that $cf$ is normalized for that
map. Since $5P$ and $4P$ avoid both its zero and its pole,
homomorphicity cancels this constant and gives
\[
 \delta(P)=\frac{f(5P)}{f(4P)}
       =\frac{-3456}{2916}
       =-\frac{32}{27}
       \equiv2^5\,3^6\pmod{(K^\times)^9}.               \tag{B-RD29.7}
\]
The last ratio is the ninth power $(-3)^9$. One can also compare
this calculation directly with the $N=9$ family and Figure~9 on
pp.26--27 of that source, specialized at $\lambda=1$.
Thus (B-RD29.1) is the Kummer extension with the specified
$q=23328$, whose valuation over $K$ is twelve. The preceding
lemma proves (B-RD29.2) with this explicit Tate parameter.

The exact invariants of $B$ are
\[
 c_4=17\,3^4,\qquad \Delta=-2\,3^9,
 \qquad j=-\frac{17^3\,3^3}{2}.                          \tag{B-RD29.8}
\]
The integral $j$ gives potential good reduction by
\cite[Section~2, p.496]{SerreTate1968Good}.
Over the completed maximal unramified extension of $K$, take
an extension of minimal good-reduction degree $m$. Its inertia
acts faithfully on the good special elliptic curve.
The displayed equation has discriminant valuation eighteen.
A change to a unit-discriminant good equation over the extension
therefore implies
\[
                            12\mid18m,
\]
so $m$ is even. The inertia image contains an element of order
two. In residue characteristic three the only nontrivial
order-two origin-preserving elliptic automorphism is $[-1]$;
it acts as $-1$ on the rational Dieudonne module and has trace
$-2$. The Tate comparison has trace $2$, by its two good graded
pieces. Hence the requested characters differ.
The inertia calculation uses the actual finite descent action
of the good model and therefore has the conventions of
Proposition~\ref{prop:B-RD4-dictionary}.

For $e\equiv2\pmod4$, take an extension of $K$ generated by a
root of $X^{e/2}-(\zeta_3-1)$. This polynomial is Eisenstein,
so its absolute ramification is $e$ and residue field remains
$\mathbf F_3$. The isomorphism (B-RD29.2) survives base change.
The raw discriminant valuation becomes $9e$, and
$12\mid9em$ still forces $m$ even. Thus the same trace
comparison holds. Adding $h-2$ identical constant etale factors
to both groups gives every $h\geq2$. Finally level nine is
$M(1)$, so its failure proves $\nu_{\min}\geq2$.
\end{proof}

\begin{remark}[Exact certificate and remaining question]
The program \path{code/verify_p3_character_level2.py} verifies all
nine subgroup points and all $81$ additions, the degree-nine
isogeny identity, the change to $B$, both discriminants and
$j$-invariants, the two divisor identities, the Kummer ratio and
its ninth-power adjustment. Its recorded exact run passed in
$4.882327191000513$ seconds; the independent parent rerun passed
in $4.293367760001274$ seconds. The program verifies finite
algebraic assertions; the Tate extension dictionary and inertia
argument are proved above.
This proves that the full-torsion bound $N_{\mathrm w}=1$ at
$e=2$ is strictly smaller than the least B.4 index. Whether
$\nu=2$ already suffices for arbitrary height-two groups at
$(p,e)=(3,2)$ remains a separate, narrower question.
\end{remark}
